% Generated draft: requires mathematical and source review.
\documentclass[11pt]{article}
\usepackage[margin=25mm]{geometry}
\usepackage{fontspec}
\setmainfont{Latin Modern Roman}
\usepackage{amsmath,amssymb,mathtools,unicode-math}
\setmathfont{Latin Modern Math}
\usepackage{enumitem,needspace,xurl,hyperref}
\setlength{\parindent}{0pt}
\setlength{\parskip}{5pt}
\setlength{\emergencystretch}{3em}
\newcommand{\N}{\mathbb{N}}
\newcommand{\Z}{\mathbb{Z}}
\newcommand{\Q}{\mathbb{Q}}
\newcommand{\R}{\mathbb{R}}
\newcommand{\C}{\mathbb{C}}
\newcommand{\F}{\mathbb{F}}
\newcommand{\A}{\mathbb{A}}
\newcommand{\PP}{\mathbb{P}}
\newcommand{\E}{\mathbb{E}}
\newcommand{\eps}{\varepsilon}
\newcommand{\dd}{\,\mathrm{d}}
\begin{document}
\section*{UnsolvedMath ID: 30000632}
\textbf{Source problem code:} catalogue entry
\subsection{Definitions and mathematical statement}
\textbf{One-sided concentration question.} Does there exist a sequence of positive numbers \(\varepsilon_n\to0\) such that for every dimension \(n\ge1\) and every isotropic origin-symmetric convex body \(K\subset\R^n\) of volume one, with isotropic constant \(\rho\),\[\operatorname{vol}_n\left\{x\in K:\frac{\|x\|_2}{\rho\sqrt n}-1\ge\varepsilon_n\right\}\le\varepsilon_n?\] The sequence must depend only on dimension, not on \(K\) or its isotropic constant; no rate of decay is prescribed. The boundary comparison \(\ge\) and the one-sided event are retained exactly.\par In the usual convention a convex body is compact convex with nonempty interior, symmetry means \(K=-K\), and isotropy at volume one means\[\int_Kx\,dx=0,\qquad \int_Kx_ix_j\,dx=\rho^2\delta_{ij}\quad(1\le i,j\le n),\] with \(\rho>0\). Thus the covariance of the uniform random point \(X\) is \(\rho^2I\), and \(\mathbb E\|X\|_2^2=n\rho^2\). The normalization \(\rho\sqrt n\) is the root-mean-square radius, not an assertion that \(\mathbb E\|X\|_2\) equals it. Since \(\operatorname{vol}_n(K)=1\), the displayed volume is also the corresponding uniform probability.\par Although the source title calls this a thin-shell hypothesis, its formula only bounds the upper radial tail. It is not replaced by \(\mathbb P(|\|X\|_2/(\rho\sqrt n)-1|\ge\varepsilon_n)\le\varepsilon_n\), a uniform variance bound, or a specified shell width. Any equivalence or strengthening to such formulations needs an additional argument and is not asserted here. No smoothness, unconditional symmetry, or fixed value of \(\rho\) is imposed.\par The supplied triage checked 2026-08-21 labels the assertion open and reports no located complete resolution, with the report as its sole source. It supplies no precise later concentration theorem or quantitative comparison. This dated status claim is preserved without independently asserting that this exact weak one-sided formulation remains unresolved today.
\subsection{Short English statement}
Seek an upper-tail concentration bound around the root-mean-square radius with one error sequence uniform over all isotropic symmetric volume-one bodies. The event is one-sided, exactly as supplied. Neither a two-sided shell statement nor a dimension-free variance bound is substituted.
\subsection{Sources}
\begin{itemize}\item Local UnsolvedMath registry, permanent ID 30000632, code catalogue entry; \emph{Thin-Shell Concentration for Isotropic Convex Bodies}, original label thin-shell \(\varepsilon\)-concentration hypothesis; Oberwolfach Reports Open Problems.\item \emph{Konvexgeometrie (2007)}, Oberwolfach Reports, DOI 10.4171/owr/2006/56, pp. 3321--3396. \url{https://doi.org/10.4171/owr/2006/56}. Proposed year and year tag: 2006; no proposer named.\item The same full report is the sole supplied literature source; no separate authoritative concentration theorem is supplied.\item Literature triage checked 2026-08-21, status \texttt{open}; record updated 2026-09-16. The \texttt{corrected\_verified} field expressly records that the one-sided event was preserved. Verify the dated unresolved-status claim independently.\end{itemize}
\end{document}
